\section{SHRINCS-SHA2-128s Specification}

\subsection{Tweak Structure and Domain Separation}
Tweaks are structured as follows to ensure uniqueness across all scheme components:

\begin{center}
\texttt{ADRS = layer || tree\_address || type || words}
\end{center}
where: 
\begin{itemize}
    \item[] \texttt{layer}: (4 bytes) layer in hypertree (\texttt{0} = bottom, \texttt{d-1} = max)
    \item[] \texttt{tree\_address}: (12 bytes) tree address within layer
    \item[] \texttt{type}: (4 bytes) address type
    \item[] \texttt{words}: (12 bytes) type-dependent 12 bytes
\end{itemize}
A total \texttt{ADRS} size is 32 bytes.

\subsubsection{Address Types}
\begin{itemize}
    \item[] \texttt{0x00}: Stateful WOTS+C chains (\texttt{SF\_WOTS\_CHAIN})
    \item[] \texttt{0x01}: Stateful tree nodes (\texttt{SF\_TREE})
    \item[] \texttt{0x02}: Stateless WOTS+C chains (\texttt{SL\_WOTS\_CHAIN})
    \item[] \texttt{0x03}: Stateless tree nodes (\texttt{SL\_TREE})
    \item[] \texttt{0x04}: FORS (\texttt{SL\_FORS})
    \item[] \texttt{0x05}: Message hash (\texttt{H\_MSG})
    \item[] \texttt{0x06}: Pseudorandom function (\texttt{PRF})
    \item[] \texttt{0x07}: Root public key (\texttt{ROOT})
\end{itemize}
Summarizing, \texttt{(0x00, 0x01)} are stateful only types, \texttt{(0x02, 0x03, 0x04)} are stateless only and \texttt{(0x05, 0x06, 0x07)} are shared types.

\subsubsection{Address Manipulation Functions}
\begin{verbatim}
setLayerAddress(ADRS, layer):
    ADRS[0:4] ← toByte(layer, 4)

setTreeAddress(ADRS, tree_addr):
    ADRS[4:16] ← toByte(tree_addr, 12)

setTypeAndClear(ADRS, type):
    ADRS[16:20] ← toByte(type, 4)
    ADRS[20:32] ← toByte(0, 12)

setKeyPairAddress(ADRS, keypair):
    ADRS[20:24] ← toByte(keypair, 4)

setChainAddress(ADRS, chain):
    ADRS[24:28] ← toByte(chain, 4)

setHashAddress(ADRS, hash):
    ADRS[28:32] ← toByte(hash, 4)

setTreeHeight(ADRS, height):
    ADRS[24:28] ← toByte(height, 4)

setTreeIndex(ADRS, index):
    ADRS[28:32] ← toByte(index, 4)
\end{verbatim}

Note that \texttt{setTreeHeight} and \texttt{setChainAddress} both write to bytes \texttt{[24:28]}. These operations are context-dependent so there mustn't be a conflict.

\oleks{TODO: type-dependent words}


\subsection{WOTS+C}
The standard Winternitz One-Time Signature (WOTS+) is the building block of most modern hash-based 
schemes. It works by encoding the message into a series of small integers and using them to iterate 
hash chains. A significant portion of a WOTS+ signature size (approx 20\%) is dedicated to a 
"Checksum," which is necessary to prevent a forgery attack where an adversary extends the hash 
chains forward.

WOTS+C (Compressed WOTS+) eliminates the transmission of this checksum entirely, achieving a smaller 
signature size. In WOTS+, the checksum $C$ is calculated from the message digest $D$.

$$C = \sum_{i=1}^{l} (w - 1 - d_i)$$

This checksum is then signed using additional chains. In WOTS+C, instead of calculating a checksum 
for a fixed message digest, the signer modifies the message digest until it possesses a specific, 
pre-determined checksum property. This is done by appending a counter \texttt{count} and a random 
salt \texttt{R} to the message before hashing. The signer increments cnt until the resulting hash digest 
$D$ satisfies a global target condition.

\subsubsection{The Target Condition}
For SHRINCS-SHA2-128s+C, we use the Winternitz parameter $\color{red}w=256$. This means each chain processes 8 
bits (1 byte) of the digest. The message digest length is 128 bits (16 bytes), so $\color{red}l = 16$. The 
values $d_i$ are treated as integers in $0\dots (w-1)$.

We define the expected value of the checksum. If the hash output is uniform, the expected value of a 
single digit is $(w-1)/2 = 127.5$. For 16 digits, the expected sum is $16 \times 127.5 = 2040$.

The condition for a valid WOTS+C message digest $D$ is:
$$\sum_{i=0}^{15}d_i = 2040$$

By enforcing this condition during signing (including grinding), the verifier can simply assume 
the checksum is 2040 without receiving it. If the reconstructed digest does not sum to 2040, the 
signature is invalid. 

\subsubsection{Parameters}

\begin{center}
\begin{tabular}{ c c }
 Parameter & SHRINCS-SHA2-128s \\ 
 \hline
 \texttt{n} & 128 bits (16 bytes) \\
 \texttt{w} & \color{red} 256 \\
 \texttt{len} & \color{red} 16 chains \\
 \texttt{Swn} & \color{red} 2040 \\
 \texttt{c} & \color{red} 4 bytes
\end{tabular}
\end{center}

\subsubsection{Base-w Representation}
\begin{verbatim}
    Algorithm: base_w(m, w, out_len)
    Input: 
        m: byte string
        w: Winternitz parameter (256)
        out_len: output length (16)
    Output: 
        b_w: array of out_len integers in [0, w-1]
    
    1. in_idx ← 0
    2. bits ← 0
    3. total ← 0
    4. b_w ← []
    5. for i from 0 to out_len - 1:
        a. if bits == 0:
            total ← m[in_idx]
            in_idx ← in_idx + 1
            bits ← 8
        b. bits ← bits - log2(w)
        c. base_w[i] ← (total >> bits) AND (w - 1)
    6. return b_w
    // for w = 256, each byte becomes one base-w digit
\end{verbatim} 

\subsubsection{Chain Function}
\begin{verbatim}
    Algorithm: chain(m, start, steps, PK.seed, ADRS)
    Input:
        m: (n/8)-byte input
        start: starting index in chain
        steps: number of iterations
        PK.seed: public seed
        ADRS: tweak
    Output:
        (n/8)-byte chain output

    1. if steps == 0:
        return m
    2. tmp ← chain(m, start, steps - 1, PK.seed, ADRS)
    3. setHashAddress(ADRS, start + steps - 1)
    4. return Tw(PK.seed, ADRS, tmp)
\end{verbatim}

\subsubsection{Key Generation}
\begin{verbatim}
    Algorithm: wots_PKgen(SK.seed, PK.seed, ADRS)
    Input:
        SK.seed: secret seed
        PK.seed: public seed
        ADRS: tweak 
    Output:
        pk: compressed WOTS+C public key ((n/8) bytes)

    1. setTypeAndClear(ADRS, WOTS_PRF)
    2. setKeyPairAddress(ADRS, ADRS.getKeyPairAddress())
    3. tmp ← []
    4. for i from 0 to l - 1:
        a. setChainAddress(ADRS, i)
        b. sk_i ← PRF(SK.seed, ADRS)
        c. setTypeAndClear(ADRS, WOTS_HASH)
        d. setKeyPairAddress(ADRS, ADRS.getKeyPairAddress())
        e. setChainAddress(ADRS, i)
        f. tmp[i] ← chain(sk_i, 0, w - 1, PK.seed, ADRS)
    5. setTypeAndClear(ADRS, WOTS_PK)
    6. setKeyPairAddress(ADRS, ADRS.getKeyPairAddress())
    7. pk ← Tw(PK.seed, ADRS, tmp[0] || ... || tmp[l-1])
    8. return pk
\end{verbatim}


\subsubsection{Messsage Digest Grinding}
\begin{verbatim}
    Algorithm: wots_grind(m, R, PK.seed, PK.root)
    Input:
        m: message to sign
        R: randomness 
        counter: counter for grinding
        PK.seed, PK.root: public key components
    Output:
        msg: array of l integers
        counter: final value used for grinding

    1. counter ← 0 
    2. for counter from 0 to 2^32:    
        a. digest ← H_msg(R || counter, PK.seed, PK.root, m)[0:n]
        b. msg ← base_w(digest, w, l)
        c. sum ← 0
        d. for i from 0 to l - 1:
            sum ← sum + msg[i]
        e. if (sum == Swn)
            return (msg, counter)
        f. counter ← counter + 1
        g. if counter == 2^32:
            abort with error // R must be resampled
\end{verbatim}
Expected iterations value is $\approx 2000\text{-}3000$ for SHRINCS-SHA2-128s with \texttt{Swn = 2040}.

\subsubsection{Signature Generation}
\begin{verbatim}
    Algorithm: wots_sign(m, SK.seed, SK.prf, PK.seed, PK.root, ADRS)
    Input:
        m: message 
        SK.seed: secret seed
        SK.prf: PRF key for randomness
        PK.seed: public seed
        PK.root: public key root
        ADRS: tweak   
    Output:
        sig: WOTS+C signature instance

    1. R ← PRF_msg(SK.prf, ADRS, m)
    2. (msg, counter) ← wots_grind(m, R, PK.seed, PK.root)
    3. sig ← R || counter
    4. setTypeAndClear(ADRS, WOTS_PRF)
    5. for i from 0 to l - 1:
        a. setChainAddress(ADRS, i)
        b. sk_i ← PRF(SK.seed, ADRS)
        c. setTypeAndClear(ADRS, WOTS_HASH)
        d. setChainAddress(ADRS, i)
        e. sig ← sig || chain(sk_i, 0, msg[i], PK.seed, ADRS)
    6. return sig
\end{verbatim}

\subsubsection{Signature Verification (Public Key Recovery)}
\begin{verbatim}
    Algorithm: wots_pkFromSig(sig, m, PK.seed, PK.root, ADRS)
    Input:
        sig: WOTS+C signature
        m: signed message
        PK.seed: public seed
        PK.root: public key root
        ADRS: tweak
    Output:
        pk: recovered public key (n bytes), or $\bot$ if invalid

    1. R ← sig[0:32]
    2. counter ← sig[32:36]
    3. (msg, valid) ← wots_digest(m, R, counter, PK.seed, PK.root) // like grinding but with known counter, should be revised
    4. if not valid:
        return $\bot$
    5. tmp ← []
    6. setTypeAndClear(ADRS, WOTS_HASH)
    7. for i from 0 to l - 1:
        a. setChainAddress(ADRS, i)
        b. sig_i ← sig[36 + i*n : 36 + (i+1)*n]
        c. tmp[i] ← chain(sig_i, msg[i], w - 1 - msg[i], PK.seed, ADRS)
    8. setTypeAndClear(ADRS, WOTS_PK)
    9. pk ← Tw(PK.seed, ADRS, tmp[0] || ... || tmp[l-1])
    10. return pk
\end{verbatim}

\subsection{UXMSS. The Stateful Tree}
Standard XMSS uses a balanced binary tree, which minimizes the maximum path length but imposes a 
logarithmic cost on every signature. For \texttt{N} signatures, every signature has size proportional to 
\texttt{log\_2(N)}. SHRINCS utilizes an Unbalanced Merkle Tree (specifically a "Right-Skewed" tree) to 
optimize for the typical Bitcoin use case: a wallet that performs very few transactions in its 
lifecycle (we are targeting \texttt{N<=10}).

In this topology the \texttt{q}-th signature requires a path of length \texttt{q+1}. This linear growth is superior 
to logarithmic growth for small \texttt{q}. For \texttt{q=0}, the path is just 16 bytes. Even for \texttt{q=10}, the path 
is $11 \times 16 = 176$ bytes, which is extremely small compared to the fixed $\approx 2.5$KB of a balanced 
tree or the $\approx 8$KB of SLH-DSA.

The rightmost leaf at maximum depth contains the hash of the stateless public key, enabling the 
fallback mechanism.

%Picture of the tree

\subsubsection{UXMSS Authentication Path}
For the \texttt{q}-th signature (\texttt{q} = count after signing), the authentication path contains \texttt{q} nodes:
\begin{center}
    \texttt{auth\_path = [node\_0, node\_1, ..., node\_{q-1}]}
\end{center}

Root computation from auth path:

\begin{verbatim}
    Algorithm: root_from_auth(pk, auth, q, PK.seed, ADRS)
    Input:
        pk: OTS public key at leaf position (q-1)
        auth: authentication path (q nodes)
        q: signature count after signing
        PK.seed: public seed
        ADRS: tweak
    Output:
        root: computed tree root

    1. node ← pk
    2. setTypeAndClear(ADRS, TREE)
    3. for i from 0 to q - 1:
        a. setTreeHeight(ADRS, i + 1)
        b. setTreeIndex(ADRS, 0)
        c. if i < q - 1:
            // auth[i] is left sibling
            node ← Tw(PK.seed, ADRS, auth[i] || node)
        d. else:
            // node is left child
            node ← Tw(PK.seed, ADRS, node || auth[i])
    4. return node
\end{verbatim}


\subsection{SLH-DSA. The Stateless Fallback}

\subsubsection{Parameters}

\begin{center}
\begin{tabular}{ c c }
 Parameter & SHRINCS-SHA2-128s \\ 
 \hline
 \texttt{h} & 63 (total tree height)\\
 \texttt{d} & 7 (number of hypertree layers) \\
 \texttt{h'} & 9 (height per layer) \\
 \texttt{a} & 12 (FORS tree height) \\
 \texttt{k} & 14 (number of FORS trees) \\
 \texttt{w} & 256 \\
 \texttt{n} & 128 \\
\end{tabular}
\end{center}


\subsection{Size}
\subsubsection{Stateful signature size}

The WOTS+C component is 292 bytes: randomness (32 bytes) + counter (4 bytes) + chain values (256 bytes).
The authentication path for signature \texttt{q} (count after signing) is \texttt{q * (n/8) = q * 16} 
bytes. The stateless public key \texttt{PK\_sl} adds \texttt{(n/8) = 16} bytes.
\begin{center}
    \texttt{sig\_sl = 292 + q * 16 + 16 = 308 + q * 16} bytes,
\end{center}
where \texttt{q} is the signature count after signing:
\begin{center}
\begin{tabular}{ c c c}
    q & Auth path & Total \\
    \hline
    1 & 16 bytes & 324 bytes \\
    2 & 32 bytes & 340 bytes \\
    3 & 48 bytes & 356 bytes \\
    10 & 160 bytes & 468 bytes
\end{tabular}
\end{center}

\subsubsection{Stateless signature size}

\subsection{SHRINCS Core Algorithms}

\subsubsection{Key Generation}

\begin{verbatim}
    Algorithm: SHRINCS.KeyGen()
    Output:
        SK: secret key
        PK: public key  
        state: initial state

    1. seed ←$ {0,1}^384
    2. SK.seed ← seed[0:16]
    3. SK.prf ← seed[16:32]
    4. PK.seed ← seed[32:48]
    5. ADRS ← new_ADRS()

    // Generate stateful public key (UXMSS root)
    6. setLayerAddress(ADRS, 0)
    7. setTreeAddress(ADRS, 0)
    8. PK_sf ← compute_initial_stateful_root(SK.seed, PK.seed, ADRS)

    // Generate stateless public key (SPHINCS+C root)
    9. PK_sl ← sphincs_root(SK.seed, PK.seed, ADRS)

    // Combined params
    10. PK.root ← H(0x08 || PK.seed || PK_sf || PK_sl)
    11. SK ← (SK.seed, SK.prf, PK.seed, PK.root, PK_sf, PK_sl)
    12. PK ← (PK.seed, PK.root)
    13. state ← { q: 0, valid: true }
    14. return (SK, PK, state)
\end{verbatim}

\subsubsection{Recovery from Seed}

\begin{verbatim}
    Algorithm: SHRINCS.Restore(seed)
    Input:
        seed: master seed
    Output:
        SK: secret key
        PK: public key
        state: state "LOST"

    1. SK.seed ← seed[0:16]
    2. SK.prf ← seed[16:32]
    3. PK.seed ← seed[32:48]
    4. ADRS ← new_ADRS()
    5. PK_sf ← compute_max_stateful_root(SK.seed, PK.seed, ADRS)
    6. PK_sl ← sphincs_root(SK.seed, PK.seed, ADRS)
    7. PK.root ← H_n(0x08 || PK.seed || PK_sf || PK_sl)
    8. SK ← (SK.seed, SK.prf, PK.seed, PK.root, PK_sf, PK_sl)
    9. PK ← (PK.seed, PK.root)
    10. state ← { q: $\bot$, valid: false } 
    11. return (SK, PK, state)
\end{verbatim}

\subsubsection{Statefull signing}

\subsubsection{Stateless signing}

\subsubsection{Statefull verification}

\subsubsection{Stateless verification}

\subsection{Data structures}

\subsubsection{Secret Key Format}

\subsubsection{Public Key Format}

\subsubsection{Stateful Signature Format}

\subsubsection{Stateless Signature Format}